% 本页为虚构示例。将院校、项目、分工、日期和验证结果换成自己的经历。
\sectionTitle{教育背景}{\faGraduationCap}
\datedsubsection{\textbf{示例大学} / \textit{微电子科学与工程} / \textit{学士}}{2022.09 -- 2026.06}
\normalline{专业方向：数字集成电路设计}
\normalline{相关课程：数字逻辑、计算机组成原理、集成电路设计基础、数字信号处理}

\sectionTitle{实习经历}{\faBriefcase}
\datedsubsection{\textbf{示例微电子公司} / 数字验证实习生}{2025.07 -- 2025.08}
\begin{itemize}
  \item \textbf{接口测试}\quad 依据 APB 寄存器表编写读写用例，覆盖复位默认值、只读位和连续访问。
  \item \textbf{缺陷复核}\quad 从失败日志提取首个读值差异，附上地址、预期值和波形，复测修订后的 RTL。
\end{itemize}

\sectionTitle{项目经历}{\faUsers}
% 先交代设计目标与个人分工，再写设计决策、验证方法和可核对的产物。
\datedsubsection{\textbf{单周期 RISC-V 处理器设计}}{2025.09 -- 2025.12}
\role{Verilog / RV32I / Verilator}{数据通路设计与指令验证}
\normalline{在教学处理器上实现算术、访存与条件分支指令，负责译码、写回选择和指令级自测。}
\begin{itemize}
  \item \textbf{数据通路}\quad 复用加法器计算访存地址，通过写回选择信号区分算术结果与加载数据。绘制控制信号表，逐项标注各类指令的寄存器读写行为。
  \item \textbf{指令验证}\quad 用短汇编程序逐条比对 PC 与寄存器写回值，检查顺序执行与分支跳转。
  \item \textbf{故障定位}\quad 从向后跳转的波形定位立即数符号扩展错误，修正后将该指令序列加入回归。
\end{itemize}

\datedsubsection{\textbf{同步 FIFO 设计与边界验证}}{2025.03 -- 2025.06}
\role{SystemVerilog / GTKWave}{RTL 实现与自检平台}
\begin{itemize}
  \item \textbf{读写控制}\quad 仅在读写被接受时更新计数器，绘制满写、空读时的接口时序图。
  \item \textbf{结果比对}\quad 用参考队列记录已接受的写入数据，逐次比对读出顺序。将空读、满写与指针回绕用例加入回归，保存预期波形供复核。
\end{itemize}

\sectionTitle{专业技能}{\faGears}
\begin{skilllist}
  \item \textbf{RTL 设计}\quad 使用 Verilog / SystemVerilog 实现组合逻辑、同步状态机与寄存器数据通路。
  \item \textbf{仿真验证}\quad 使用 Verilator 运行指令测试，通过 GTKWave 检查时钟边沿与控制信号。
  \item \textbf{开发环境}\quad 在 Linux 下用 Make 组织仿真任务，用 Git 记录 RTL 与测试用例的版本。
\end{skilllist}

% 可选模块：按目标岗位决定是否保留课程设计；也可替换为实习、科研或作品。
\sectionTitle{课程设计}{\faFlask}
\datedsubsection{\textbf{含噪信号的 FIR 低通滤波}}{2024.09 -- 2024.12}
\role{MATLAB / FIR}{滤波器设计与频域分析}
\begin{itemize}
  \item 根据测试信号和干扰频率设置通带与阻带，比较不同阶数下的过渡带宽度与群时延。
  \item 将输入、输出频谱置于同一坐标尺度，整理参数表和绘图脚本，供同组同学复现实验。
\end{itemize}

\sectionTitle{荣誉奖项}{\faAward}
\datedline{校级电子设计竞赛二等奖 / 数字系统设计组}{2025.05}
\datedline{学业优秀奖学金 / 示例大学}{2024.10}

% 可选替换模块；以下均为虚构填写示例。
% \sectionTitle{技术作品}{\faFolderOpen}
% \datedsubsection{\textbf{RTL 仿真用例集}}{2025.12}
% \normalline{整理 FIFO 与译码器的最小测试用例，附运行命令、预期输出和失败波形。}
% \sectionTitle{科研经历}{\faFlask}
% \datedsubsection{\textbf{低功耗计数器的翻转活动分析} / 校内课题组}{2025.03 -- 2025.06}
% \normalline{在同一组输入激励下记录两种编码的信号翻转次数，整理测试条件与对照结果。}
